\documentclass[12pt]{article}
\usepackage[a4paper, margin=1in]{geometry}
\usepackage{fontspec}
\usepackage{polyglossia}
\usepackage{amsmath}
\usepackage{amssymb}
\usepackage{graphicx}
\usepackage{listings}
\usepackage{color}
\usepackage{enumitem}

\setmainlanguage{thai}
\setmainfont{Angsana New}

\newfontfamily\thaifont{TH SarabunPSK}[Script=Thai]
\newfontfamily\thaifonttt{Courier New} % หรือฟอนต์ Monospace อื่นๆ สำหรับโค้ด

\definecolor{codegray}{rgb}{0.5,0.5,0.5}
\definecolor{codepurple}{rgb}{0.58,0,0.82}
\definecolor{backcolour}{rgb}{0.95,0.95,0.92}

\lstset{
    backgroundcolor=\color{backcolour},   
    commentstyle=\color{codegray},
    keywordstyle=\color{blue},
    numberstyle=\tiny\color{codegray},
    stringstyle=\color{codepurple},
    basicstyle=\ttfamily\small,
    breakatwhitespace=false,         
    breaklines=true,                 
    captionpos=b,                    
    keepspaces=true,                 
    numbers=left,                    
    numbersep=5pt,                  
    showspaces=false,                
    showstringspaces=false,
    showtabs=false,                  
    tabsize=2
}

\title{รายงานสรุปข้อสอบและเฉลยรายวิชาด้าน Data Science และ Machine Learning}
\author{Gemini CLI Autonomous Agent}
\date{\today}

\begin{document}

\maketitle
\newpage

\section{Learn Python: Language Basics and Fundamental Data Processing}

\subsection*{ข้อที่ 1}
\textbf{โจทย์:} Dynamically-typed language หมายความว่าอะไร\\
\textbf{คำตอบ:} ประเภทของตัวแปรถูกตัดสินตอน Run-Time จึงไม่ต้องกำหนดประเภทของตัวแปรก่อนใช้งาน\\
\textbf{คำอธิบาย:} ภาษาโปรแกรมที่มีลักษณะเป็น Dynamically-typed คือภาษาที่ "ประเภทของข้อมูล" ของตัวแปรจะถูกผูกติดอยู่กับค่า (Value) ที่เก็บอยู่ในตัวแปรในขณะนั้น การตัดสินประเภทข้อมูลจะเกิดขึ้นในขณะที่โปรแกรมกำลังทำงาน (Run-time)

\subsection*{ข้อที่ 2}
\textbf{โจทย์:} อะไรเป็นหัวใจหลักที่ทำให้ Jupyter notebook เข้าใจภาษาโปรแกรมภาษาต่าง ๆ\\
\textbf{คำตอบ:} Kernel\\
\textbf{คำอธิบาย:} สถาปัตยกรรมของ Jupyter แยกส่วนแสดงผล (Notebook) ออกจากส่วนประมวลผล (Kernel) การรองรับหลายภาษาทำได้โดยการติดตั้ง Kernel เฉพาะของภาษานั้นๆ

\subsection*{ข้อที่ 3}
\textbf{โจทย์:} ถ้าสมุดโน้ต ณ ตอนนี้มี 5 cell และ cell ปัจจุบันอยู่ที่ cell รองสุดท้าย หลังจากที่รัน cell ปัจจุบันด้วยการกดปุ่ม Shift+Enter สมุดโน้ตจะมีกี่ cell\\
\textbf{คำตอบ:} 5 cell\\
\textbf{คำอธิบาย:} เมื่อรัน Cell รองสุดท้าย (ลำดับที่ 4 จาก 5) ด้วย Shift+Enter ระบบจะเลื่อน Focus ไปยัง Cell ถัดไป (ลำดับที่ 5) โดยไม่มีการสร้าง Cell ใหม่ จำนวน Cell รวมจึงยังคงเป็น 5

\subsection*{ข้อที่ 4}
\textbf{โจทย์:} หากต้องการแสดง block ของโค้ดใน Markdown cell เราสามารถทำได้อย่างไร\\
\textbf{คำตอบ:} ใส่ space จำนวนเท่ากันแต่ต้องอย่างน้อย 4 ตัวในส่วนต้นของบรรทัดของโค้ดที่อยู่ใน Block เดียวกัน\\
\textbf{คำอธิบาย:} การเว้นวรรค 4 ช่องหรือ 1 Tab เป็นมาตรฐานดั้งเดิม (Indented Code Block) ของ Markdown ในการสร้างส่วนแสดงผลโค้ด

\subsection*{ข้อที่ 5}
\textbf{โจทย์:} ข้อใดเป็นผลลัพธ์ของโปรแกรมต่อไปนี้
\begin{lstlisting}[language=Python]
x = 2
y = 4
y /= x
print(y)
\end{lstlisting}
\textbf{คำตอบ:} 2.0\\
\textbf{คำอธิบาย:} ตัวดำเนินการ \texttt{/} ใน Python 3 จะให้ผลลัพธ์เป็นทศนิยม (Float) เสมอ แม้การหารจะลงตัวก็ตาม

\subsection*{ข้อที่ 6}
\textbf{โจทย์:} expression 4 + - 1 * 4 / 2 ** 2 มีค่าเท่ากับเท่าใด\\
\textbf{คำตอบ:} 3.0\\
\textbf{คำอธิบาย:} ลำดับความสำคัญเริ่มจากเลขยกกำลัง (\texttt{2**2=4}) -> คูณ/หาร (\texttt{-1*4/4 = -1.0}) -> บวก (\texttt{4 + (-1.0) = 3.0})

\subsection*{ข้อที่ 7}
\textbf{โจทย์:} ให้ function funcA() มี definition ส่วนหัวเป็น funcA( *arg1 , arg2 = True ) ในการเรียกใช้ funcA() จะต้องมีการใส่ argument จำนวนกี่ค่า\\
\textbf{คำตอบ:} 0 ค่า\\
\textbf{คำอธิบาย:} เนื่องจาก \texttt{*arg1} เป็นการรับ Argument แบบไม่จำกัดจำนวน (รวมถึง 0 ค่า) และ \texttt{arg2} มีค่าเริ่มต้น (Default value) จึงสามารถเรียกใช้ฟังก์ชันโดยไม่ส่ง Argument ใดๆ ได้เลย

\subsection*{ข้อที่ 8}
\textbf{โจทย์:} ข้อใดเกิด implicit data type conversion และผลลัพธ์มี data type ใด\\
\textbf{คำตอบ:} 3 + 4.0 ผลลัพธ์เป็น float\\
\textbf{คำอธิบาย:} เมื่อมีการดำเนินการระหว่าง Integer และ Float ระบบจะแปลง Integer เป็น Float โดยอัตโนมัติ (Type Promotion) เพื่อความแม่นยำ

\subsection*{ข้อที่ 9}
\textbf{โจทย์:} ผลลัพธ์เมื่อรัน print(type("34" + 3)) คือข้อใด\\
\textbf{คำตอบ:} เกิด error\\
\textbf{คำอธิบาย:} Python เป็น Strong Typing ไม่ยอมให้เชื่อม (Concatenate) String เข้ากับ Integer โดยตรง จะเกิด \texttt{TypeError}

\subsection*{ข้อที่ 10}
\textbf{โจทย์:} หลังจากรันโปรแกรมต่อไปนี้ ค่า x จะมีค่าเท่ากับปริมาณใด
\begin{lstlisting}[language=Python]
x = 0
for i in range(1,1000,2):
    if i%2 == 0:
        x += i
\end{lstlisting}
\textbf{คำตอบ:} 0\\
\textbf{คำอธิบาย:} \texttt{range(1, 1000, 2)} ให้เลขคี่ทั้งหมด ซึ่งไม่มีทางที่ \texttt{i\%2 == 0} จะเป็นจริง เงื่อนไขการบวกจึงไม่เคยเกิดขึ้น ค่า x จึงยังคงเป็น 0

\subsection*{ข้อที่ 11}
\textbf{โจทย์:} โปรแกรมต่อไปนี้ให้ผลลัพธ์อะไร
\begin{lstlisting}[language=Python]
x = 1
y = 1000
while x <= y:
    if abs(x-y) > 500:
        y = 100
        continue
    elif abs(x-y) > 10:
        y = 5
        continue
    else:
        y -= 1
    x += 1
print(x)
\end{lstlisting}
\textbf{คำตอบ:} 4\\
\textbf{คำอธิบาย:} โปรแกรมวนลูปโดยปรับค่า y จาก 1000 -> 100 -> 5 จากนั้นจึงเข้าสู่เงื่อนไข \texttt{else} เพื่อเพิ่มค่า x จนกระทั่ง \texttt{x > y} ผลสุดท้าย x = 4

\subsection*{ข้อที่ 12}
\textbf{โจทย์:} ข้อใดถูกต้อง\\
\textbf{คำตอบ:} Python list มี items ที่เป็นข้อมูลประเภท list ปนกับข้อมูลประเภทอื่นได้\\
\textbf{คำอธิบาย:} Python list เป็น Heterogeneous collection ที่สามารถเก็บข้อมูลต่างชนิดกันรวมถึง List ซ้อนกันได้

\subsection*{ข้อที่ 13}
\textbf{โจทย์:} ข้อใดถูกสำหรับ Python Set\\
\textbf{คำตอบ:} เป็น unordered data collection\\
\textbf{คำอธิบาย:} Set ใน Python ไม่มีลำดับ (Unordered) และไม่สามารถเข้าถึงผ่าน Index ได้ รวมถึงเก็บค่าที่ไม่ซ้ำกันเท่านั้น

\subsection*{ข้อที่ 14}
\textbf{โจทย์:} ข้อใดถูกต้อง ถ้ารันคำสั่ง t1 = tuple('1234')\\
\textbf{คำตอบ:} t1 จะมี 4 items ที่มีค่า '1', '2', '3' และ '4'\\
\textbf{คำอธิบาย:} ฟังก์ชัน \texttt{tuple()} จะแตกสมาชิกจาก Iterable (ในที่นี้คือ String) ออกมาทีละตัวละครสร้างเป็น Tuple

\subsection*{ข้อที่ 15}
\textbf{โจทย์:} ให้ product เป็น dictionary ที่มีค่า {'id': 'P123', 'name': 'watch', 'price': 5000} ถ้าต้องการเพิ่ม 'color' เป็นอีก key หนึ่ง ที่มีค่า 'silver' ต้องเขียน code อย่างไร\\
\textbf{คำตอบ:} product['color'] = 'silver'\\
\textbf{คำอธิบาย:} การเพิ่มหรืออัปเดต Key ใน Dictionary ทำได้โดยการกำหนดค่าผ่านวงเล็บก้ามปูโดยตรง

\subsection*{ข้อที่ 16}
\textbf{โจทย์:} ข้อใดไม่ใช่ประโยชน์ของ Functions\\
\textbf{คำตอบ:} Availability\\
\textbf{คำอธิบาย:} Availability (ความพร้อมใช้งาน) เป็นคุณสมบัติของระบบงานในระดับโครงสร้างพื้นฐาน ไม่ใช่ประโยชน์โดยตรงจากการจัดโครงสร้างโค้ดด้วยฟังก์ชัน

\subsection*{ข้อที่ 17}
\textbf{โจทย์:} จากฟังก์ชั่น signature: \texttt{def create\_customer(first\_name, last\_name, age, is\_new\_customer=True, **kwargs):}\\
ข้อใดทำให้เกิด error จากการเรียกใช้\\
\textbf{คำตอบ:} create\_customer('David', last\_name='Berry', title='Mr.', occupation='programmer')\\
\textbf{คำอธิบาย:} ขาด Argument ที่จำเป็น (Required Positional Argument) คือ \texttt{age}

\subsection*{ข้อที่ 18}
\textbf{โจทย์:} ข้อใดคือประโยชน์ของการเปิดไฟล์ด้วย keyword with\\
\textbf{คำตอบ:} ไม่ต้องเรียกฟังก์ชั่น file close() เพื่อปิดไฟล์ที่เปิดไว้ หลังจากออกจาก block ของ with แล้ว\\
\textbf{คำอธิบาย:} \texttt{with} ทำหน้าที่เป็น Context Manager ที่จะเรียกใช้เมธอดปิดไฟล์ให้อัตโนมัติเมื่อออกจากบล็อก

\subsection*{ข้อที่ 19}
\textbf{โจทย์:} โปรแกรมต่อไปนี้ให้ผลลัพธ์อะไร
\begin{lstlisting}[language=Python]
coor = [(1,2),(2,4),(3,6),(4,9)]
print(list(map(lambda i: i[1]==2*i[0] ,coor)))
\end{lstlisting}
\textbf{คำตอบ:} [True, True, True, False]\\
\textbf{คำอธิบาย:} ตรวจสอบว่าสมาชิกตัวที่ 2 ใน Tuple เท่ากับ 2 เท่าของตัวแรกหรือไม่ (1,2)->T, (2,4)->T, (3,6)->T, (4,9)->F

\subsection*{ข้อที่ 20}
\textbf{โจทย์:} ข้อใดคือลำดับขั้น scope ของ modules ใน Python จากใหญ่ไปเล็กที่ถูกต้อง\\
\textbf{คำตอบ:} Package -> Python File -> Function\\
\textbf{คำอธิบาย:} Package บรรจุหลายไฟล์ (Modules), ไฟล์บรรจุคลาสและฟังก์ชันตามลำดับโครงสร้าง

\newpage
\section{รู้จักคอมพิวเตอร์และระบบดิจิทัล}

\subsection*{ข้อที่ 1}
\textbf{โจทย์:} ตัวเลือกใดแสดงแนวคิดของการทำงานโดยคอมพิวเตอร์เชิงไฟฟ้า\\
\textbf{คำตอบ:} หลักการเปิด/ปิดสวิตช์ แทนค่า 0 หรือ 1 ในเลขฐานสอง\\
\textbf{คำอธิบาย:} ทรานซิสเตอร์ทำหน้าที่เป็นสวิตช์ที่มีสถานะ เปิด (1) และ ปิด (0) ซึ่งเป็นพื้นฐานของระบบดิจิทัล

\subsection*{ข้อที่ 2}
\textbf{โจทย์:} ทรานซิสเตอร์ถูกใช้อย่างไรในระบบคอมพิวเตอร์\\
\textbf{คำตอบ:} ใช้แทนสวิตช์เชิงไฟฟ้า\\
\textbf{คำอธิบาย:} ทรานซิสเตอร์ทำหน้าที่ควบคุมการไหลของกระแสไฟฟ้าเพื่อแทนค่าทางตรรกศาสตร์ (Boolean Logic)

\subsection*{ข้อที่ 3}
\textbf{โจทย์:} กำหนดข้อมูลเลขฐานสอง ขนาด 8 หลักเป็น 10010000 ข้อใดถูกต้อง\\
\textbf{คำตอบ:} ถูกทุกตัวเลือก\\
\textbf{คำอธิบาย:} (1) บิตซ้ายสุดเป็น 1 คือค่าลบ (2) ไม่คิดเครื่องหมายได้ 144 (3) 8 บิตแทนได้ 256 รูปแบบ

\subsection*{ข้อที่ 4}
\textbf{โจทย์:} หากกำหนด A = 0 และ B = 1 ตัวเลือกใดเป็นผลลัพธ์ของวงจรต่อไปนี้\\
\textbf{ลักษณะรูปภาพ:} ภาพแสดงวงจรลอจิกเกตที่ประกอบด้วย AND gate และ NOT gate ต่อร่วมกัน โดยรับอินพุต A และ B ผ่านเกตต่างๆ จนถึงเอาต์พุตสุดท้าย\\
\textbf{คำตอบ:} 1\\
\textbf{คำอธิบาย:} จากการไล่ลอจิก: NOT(A AND A) = 1, NOT(B AND B) = 0, (1 AND 0) = 0, และ NOT ตัวสุดท้ายกลับค่า 0 เป็น 1

\subsection*{ข้อที่ 5}
\textbf{โจทย์:} หากกำหนด A = 0100 และ B = 0010 ตัวเลือกใดเป็นผลลัพธ์ของ A AND B\\
\textbf{คำตอบ:} 0000\\
\textbf{คำอธิบาย:} การทำ Bitwise AND ทีละตำแหน่ง: (0\&0), (1\&0), (0\&1), (0\&0) ผลรวมเป็น 0000

\subsection*{ข้อที่ 6}
\textbf{โจทย์:} จากตัวเลือกต่อไปนี้ หน่วยความจำแบบใดทำงานได้เร็วที่สุด\\
\textbf{คำตอบ:} CPU (Register)\\
\textbf{คำอธิบาย:} Register อยู่ภายใน CPU โดยตรง มีความเร็วสูงสุดเนื่องจากไม่มีความหน่วงจากการรับส่งผ่านบัสภายนอก

\subsection*{ข้อที่ 7}
\textbf{โจทย์:} ตัวเลือกใดเป็นส่วนคำนวณของคอมพิวเตอร์\\
\textbf{คำตอบ:} CPU\\
\textbf{คำอธิบาย:} CPU (Central Processing Unit) มีหน่วย ALU ที่ทำหน้าที่คำนวณทางคณิตศาสตร์และตรรกศาสตร์

\subsection*{ข้อที่ 8}
\textbf{โจทย์:} การทำงานร่วมกันของอุปกรณ์ฮาร์ดแวร์ต่างๆ เกิดขึ้นได้อย่างไร\\
\textbf{คำตอบ:} ระบบปฏิบัติการดำเนินการควบคุมและประสานงานการทำงานส่วนต่างๆ\\
\textbf{คำอธิบาย:} OS ทำหน้าที่เป็นตัวกลางบริหารจัดการทรัพยากร (CPU, RAM, Storage) ให้ทำงานสอดประสานกัน

\subsection*{ข้อที่ 9}
\textbf{โจทย์:} ภาษาแอสเซมบลีคืออะไร\\
\textbf{คำตอบ:} ภาษาโปรแกรมซึ่งเทียบเท่ากับภาษาเครื่อง แต่เขียนให้อยู่ในรูปแบบที่อ่านได้ง่ายขึ้น\\
\textbf{คำอธิบาย:} เป็นภาษาระดับต่ำที่ใช้ตัวอักษรย่อ (Mnemonics) แทนรหัสเลขฐานสองของภาษาเครื่อง

\subsection*{ข้อที่ 10}
\textbf{โจทย์:} การมี core หรือแกนประมวลผล มากขึ้น มีประโยชน์อย่างไร\\
\textbf{คำตอบ:} ถูกทุกตัวเลือก\\
\textbf{คำอธิบาย:} ช่วยในการทำ Multitasking, ขนานการประมวลผล (Parallelism) และรองรับผู้ใช้ได้มากขึ้น

\subsection*{ข้อที่ 11}
\textbf{โจทย์:} ตัวเลือกใดเป็นหน้าที่ของระบบปฏิบัติการ\\
\textbf{คำตอบ:} ถูกทุกตัวเลือก\\
\textbf{คำอธิบาย:} ครอบคลุมการจัดสรรทรัพยากร, การควบคุมสิทธิการใช้งาน และการจัดการระบบเครือข่าย

\subsection*{ข้อที่ 12}
\textbf{โจทย์:} เหตุใดไฟล์ภาพ JPG จากมือถือจึงเปิดบนคอมพิวเตอร์ได้แม้สถาปัตยกรรมต่างกัน\\
\textbf{คำตอบ:} เพราะโปรแกรมอ่านไฟล์เข้าใจข้อตกลงการแทนข้อมูลของไฟล์\\
\textbf{คำอธิบาย:} เกิดจากมาตรฐานรูปแบบไฟล์ (File Format Standards) ที่เป็นข้อตกลงร่วมกันในการเก็บข้อมูล

\subsection*{ข้อที่ 13}
\textbf{โจทย์:} ตัวเลือกใดถูกต้องเกี่ยวกับการเก็บข้อมูลภาพ\\
\textbf{คำตอบ:} ถูกทุกตัวเลือก\\
\textbf{คำอธิบาย:} ครอบคลุมทั้งการเก็บองค์ประกอบสี (RGB), การแบ่งเป็นจุด (Bitmap/Pixel) และการใช้สมการ (Vector)

\subsection*{ข้อที่ 14}
\textbf{โจทย์:} หากข้อมูลทุกรูปแบบเก็บเป็นเลขฐานสอง ข้อใดถูกต้อง\\
\textbf{คำตอบ:} ถูกทุกตัวเลือก\\
\textbf{คำอธิบาย:} ข้อมูลดิจิทัลคือชุดตัวเลข ความหมายขึ้นอยู่กับการตีความของโปรแกรมที่นำไปเปิด

\subsection*{ข้อที่ 15}
\textbf{โจทย์:} หากกำหนด A = 0 และ B = 1 ตัวเลือกใดเป็นผลลัพธ์ของวงจรต่อไปนี้\\
\textbf{ลักษณะรูปภาพ:} ภาพแสดงสัญลักษณ์ของ XOR Gate (Exclusive OR Gate) ที่รับอินพุต A และ B\\
\textbf{คำตอบ:} 1\\
\textbf{คำอธิบาย:} XOR Gate จะให้ผลลัพธ์เป็น 1 เมื่ออินพุตต่างกัน (0 และ 1)

\newpage
\section{รู้จักปัญญาประดิษฐ์ และการเรียนรู้ของเครื่อง}

\subsection*{ข้อที่ 1}
\textbf{โจทย์:} ข้อใดเป็นคำที่ตีความได้แตกต่างกันมากที่สุดในคำนิยามของปัญญาประดิษฐ์\\
\textbf{คำตอบ:} ความฉลาด\\
\textbf{คำอธิบาย:} "ความฉลาด" เป็นนามธรรมที่ไม่มีนิยามสากลและมีการตีความหลากหลายในเชิงจิตวิทยาและวิทยาการคอมพิวเตอร์

\subsection*{ข้อที่ 2}
\textbf{โจทย์:} การทดสอบทัวริง (Turing Test) เกี่ยวข้องกับข้อใดมากที่สุด\\
\textbf{คำตอบ:} การเลียนแบบมนุษย์ด้านการสื่อสาร\\
\textbf{คำอธิบาย:} เป็นการทดสอบความสามารถของเครื่องจักรในการสื่อสารผ่านตัวอักษรให้แนบเนียนจนมนุษย์แยกไม่ออก

\subsection*{ข้อที่ 3}
\textbf{โจทย์:} ข้อได้เปรียบของระบบปัญญาประดิษฐ์ที่มีมากกว่ามนุษย์คือข้อใด\\
\textbf{คำตอบ:} คำนวณทางคณิตศาสตร์ได้ดีกว่ามนุษย์\\
\textbf{คำอธิบาย:} คอมพิวเตอร์ประมวลผลตัวเลขมหาศาลได้แม่นยำและรวดเร็วกว่ามนุษย์ที่จำกัดด้วยความเหนื่อยล้า

\subsection*{ข้อที่ 4}
\textbf{โจทย์:} ข้อใดไม่ใช่ตัวอย่างของการประยุกต์ใช้งานทางปัญญาประดิษฐ์\\
\textbf{คำตอบ:} การค้นหาประวัติคนไข้จากรหัสแท่ง (Barcode)\\
\textbf{คำอธิบาย:} การสแกนบาร์โค้ดเป็นการเข้าถึงข้อมูลจากฐานข้อมูลโดยตรง ไม่มีการเรียนรู้หรือการตัดสินใจที่ซับซ้อนแบบ AI

\subsection*{ข้อที่ 5}
\textbf{โจทย์:} การหาวิธีการเลื่อนแผ่นในเกม 8-Puzzle ไม่เกี่ยวข้องกับข้อใด\\
\textbf{คำตอบ:} การเข้าใจภาษาธรรมชาติ\\
\textbf{คำอธิบาย:} เกม 8-Puzzle แก้ไขด้วยตรรกศาสตร์และการค้นหาในพื้นที่สถานะ (State Space) ไม่เกี่ยวกับการประมวลผลภาษา

\subsection*{ข้อที่ 6}
\textbf{โจทย์:} ข้อใดคือปัญหาของการใช้การค้นหาแบบ Depth First Search (DFS)\\
\textbf{คำตอบ:} หน่วยความจำที่ต้องใช้\\
\textbf{คำอธิบาย:} DFS ต้องเก็บเส้นทางใน Stack หากกราฟลึกมากอาจเกิดปัญหาหน่วยความจำเต็มหรือ Stack Overflow

\subsection*{ข้อที่ 7}
\textbf{โจทย์:} ข้อใดเป็นข้อดีของการใช้ต้นไม้ตัดสินใจ (Decision Tree)\\
\textbf{คำตอบ:} มนุษย์สามารถเข้าใจผลลัพธ์จากโมเดลได้ง่ายกว่าการเรียนรู้ในแบบอื่น ๆ\\
\textbf{คำอธิบาย:} เป็นโมเดลแบบกล่องขาว (White-box) ที่แสดงขั้นตอนการตัดสินใจในรูปแบบ Flowchart ที่มนุษย์อ่านได้

\subsection*{ข้อที่ 8}
\textbf{โจทย์:} ในการจำแนกรูปภาพ นิยมใช้การเรียนรู้แบบใด\\
\textbf{คำตอบ:} นิวรอลเน็ตเวิร์กแบบคอนโวลูชัน (Convolutional Neural Network - CNN)\\
\textbf{คำอธิบาย:} CNN ถูกออกแบบมาเพื่อตรวจจับคุณลักษณะเชิงพื้นที่ (Spatial Features) ในรูปภาพได้โดยอัตโนมัติ

\subsection*{ข้อที่ 9}
\textbf{โจทย์:} ข้อใดคือหัวใจสำคัญของกระบวนการการเรียนรู้ของเครื่อง\\
\textbf{คำตอบ:} ลดข้อผิดพลาด (Error) บนตัวอย่างที่ใช้เรียนรู้\\
\textbf{คำอธิบาย:} การเรียนรู้คือการปรับพารามิเตอร์เพื่อทำให้ความแตกต่างระหว่างค่าทำนายและค่าจริง (Loss/Error) ต่ำที่สุด

\subsection*{ข้อที่ 10}
\textbf{โจทย์:} อัลกอริทึมใดเหมาะกับการพยากรณ์แนวโน้มค่าต่อเนื่อง เช่น ราคาหุ้น\\
\textbf{คำตอบ:} การวิเคราะห์การถดถอยเชิงเส้น (Linear Regression)\\
\textbf{คำอธิบาย:} Regression ออกแบบมาเพื่อทำนายค่าตัวเลขต่อเนื่องและระบุทิศทางความสัมพันธ์ของข้อมูล

\newpage
\section{Machine Learning Series: Supervised Learning}

\subsection*{ข้อที่ 1}
\textbf{โจทย์:} ข้อใดกล่าวผิดเกี่ยวกับ Loss function\\
\textbf{คำตอบ:} ค่า Loss เยอะๆ ยิ่งดี\\
\textbf{คำอธิบาย:} Loss Function วัดความผิดพลาด เป้าหมายคือการทำให้ค่านี้เหลือน้อยที่สุดเพื่อให้โมเดลแม่นยำที่สุด

\subsection*{ข้อที่ 2}
\textbf{โจทย์:} ข้อใดกล่าวผิดเกี่ยวกับ Hyper-parameters\\
\textbf{คำตอบ:} เป็นค่าที่มีการเปลี่ยนแปลงด้วย Learning algorithm\\
\textbf{คำอธิบาย:} Hyper-parameters ถูกกำหนดโดยผู้ใช้ก่อนเริ่มเทรน ส่วนค่าที่ปรับเปลี่ยนระหว่างเทรนคือ Model Parameters

\subsection*{ข้อที่ 3}
\textbf{โจทย์:} หากต้องการวัด Unbiased generalisation performance ต้องแยก Test set ออกจากส่วนใด\\
\textbf{คำตอบ:} ถูกทุกข้อ (Training, Model selection, Hyper-parameter tuning)\\
\textbf{คำอธิบาย:} ต้องเก็บ Test set แยกไว้ต่างหากเพื่อป้องกัน Data Leakage และใช้ประเมินผลเพียงครั้งเดียวในตอนท้าย

\subsection*{ข้อที่ 4}
\textbf{โจทย์:} ข้อใดคือคำจำกัดความของ Overfitting\\
\textbf{คำตอบ:} เมื่อเพิ่ม Model complexity แล้ว Training error ลดลง แต่ Test error เพิ่มขึ้น\\
\textbf{คำอธิบาย:} สภาวะที่โมเดลจดจำ Noise ใน Training data มากเกินไปจนไม่สามารถทำนายข้อมูลใหม่ได้ดี

\subsection*{ข้อที่ 5}
\textbf{โจทย์:} การแก้ปัญหา Overfitting ทำได้โดยวิธีใด\\
\textbf{คำตอบ:} ถูกทุกข้อ (Regularisation, ลด Degree of polynomial, เพิ่ม Training samples)\\
\textbf{คำอธิบาย:} วิธีการเหล่านี้มุ่งเน้นการลดความซับซ้อนของโมเดลหรือเพิ่มข้อมูลเพื่อให้โมเดลเรียนรู้ Pattern จริง

\subsection*{ข้อที่ 6}
\textbf{โจทย์:} Error ประเภทใดไม่ขึ้นอยู่กับ Model complexity หรือโมเดลที่เลือกใช้\\
\textbf{คำตอบ:} Irreducible error\\
\textbf{คำอธิบาย:} เป็น Noise ที่เกิดจากธรรมชาติของข้อมูล ไม่สามารถลดทอนได้ไม่ว่าจะใช้โมเดลเก่งเพียงใด

\subsection*{ข้อที่ 7}
\textbf{โจทย์:} ค่าของ Normalized root mean squared error (NRMSE) อยู่ในช่วงใด\\
\textbf{คำตอบ:} [0, $\infty$)\\
\textbf{คำอธิบาย:} RMSE เป็นบวกเสมอ และการ Normalization ด้วยค่าบวกทำให้ผลลัพธ์เริ่มที่ 0 และไม่มีขีดจำกัดบน

\subsection*{ข้อที่ 8}
\textbf{โจทย์:} การจัดความพอใจเป็น positive หรือ negative เป็น Classification ประเภทใด\\
\textbf{คำตอบ:} Binary classification\\
\textbf{คำอธิบาย:} เป็นการจัดกลุ่มข้อมูลออกเป็น 2 คลาสที่เป็นคู่ตรงข้ามกัน

\subsection*{ข้อที่ 9}
\textbf{โจทย์:} ข้อใดคือชื่อของ Binary classifier\\
\textbf{คำตอบ:} Logistic regression\\
\textbf{คำอธิบาย:} แม้ชื่อจะลงท้ายด้วย Regression แต่ทำงานโดยใช้ Sigmoid ฟังก์ชันเพื่อจัดกลุ่มข้อมูลเป็น 2 ประเภท

\subsection*{ข้อที่ 10}
\textbf{โจทย์:} ข้อใดคือ Loss function ที่เหมาะสมในการทำ Binary classification\\
\textbf{คำตอบ:} Cross entropy loss\\
\textbf{คำอธิบาย:} ถูกออกแบบมาเพื่อวัดความแตกต่างระหว่างความน่าจะเป็น (0-1) และใช้หลักการ Maximum Likelihood

\subsection*{ข้อที่ 11}
\textbf{โจทย์:} ถ้าต้องการทำ Multi-class classification ที่สามารถได้ 4 คลาส จะต้องสร้าง Binary classifier อย่างน้อยที่สุดกี่อัน โดยที่คำนึงถึง Unknown class ด้วย\\
\textbf{คำตอบ:} 4\\
\textbf{คำอธิบาย:} ใช้เทคนิค One-vs-Rest (OvR) โดยสร้าง 1 ตัวต่อ 1 คลาส หากทั้ง 4 ตัวตอบ No ทั้งหมด จะสรุปได้ว่าเป็น Unknown

\subsection*{ข้อที่ 12}
\textbf{โจทย์:} ตัวเลือกใดไม่ควรใช้เป็น Evaluation metrics สำหรับปัญหา Classification ได้\\
\textbf{คำตอบ:} Mean Squared Error (MSE)\\
\textbf{คำอธิบาย:} MSE ออกแบบมาสำหรับงาน Regression ที่ต้องการวัดระยะห่างเชิงตัวเลข ไม่เหมาะกับงานจัดหมวดหมู่ที่ระยะห่างระหว่างคลาสไม่มีความหมายเชิงเปรียบเทียบ

\subsection*{ข้อที่ 13}
\textbf{โจทย์:} ข้อใดเรียงลำดับขั้นตอนในการ Machine learning project ไม่เหมาะสม\\
\textbf{คำตอบ:} Over sampling -> Split train-test -> Train model -> Evaluate model\\
\textbf{คำอธิบาย:} การทำ Over sampling ต้องทำหลังจากการแบ่งข้อมูล (Split) แล้วเท่านั้น เพื่อป้องกัน Data Leakage

\subsection*{ข้อที่ 14}
\textbf{โจทย์:} Output activation function ควรใช้สำหรับ Binary classification\\
\textbf{คำตอบ:} Sigmoid\\
\textbf{คำอธิบาย:} Sigmoid บีบอัดค่าให้อยู่ในช่วง (0, 1) ซึ่งสอดคล้องกับคุณสมบัติของความน่าจะเป็นในงาน Binary classification

\subsection*{ข้อที่ 15}
\textbf{โจทย์:} ตัวเลือกใดเป็นการทำ Regularization ใน SVM\\
\textbf{คำตอบ:} เพิ่มค่า Slack variable (หรือลดค่า C)\\
\textbf{คำอธิบาย:} การยอมให้มี Slack variable มากขึ้น (ผ่านการลด C) ทำให้ Margin กว้างขึ้นและโมเดลยืดหยุ่นขึ้น ช่วยลด Overfitting

\subsection*{ข้อที่ 16}
\textbf{โจทย์:} ข้อใดกล่าวถูกต้องเกี่ยวกับ Ensemble Learning\\
\textbf{คำตอบ:} ถูกทุกตัวเลือก (Boosting เป็น Sequential, Stacking มี Meta-learner, Bagging เป็น Parallel)\\
\textbf{คำอธิบาย:} เป็นลักษณะเฉพาะของการทำงานในเทคนิค Ensemble แต่ละประเภท

\subsection*{ข้อที่ 17}
\textbf{โจทย์:} ตัวเลือกใดไม่ใช่จุดประสงค์ของการทำ Feature engineering\\
\textbf{คำตอบ:} เพื่อทำให้ความสัมพันธ์ระหว่าง Features กับ Target มีความซับซ้อนลดลง\\
\textbf{คำอธิบาย:} ปกติ Feature Engineering มักใช้เพื่อดึงความสัมพันธ์ที่ซ่อนอยู่ (อาจซับซ้อนขึ้นแต่ชัดเจนขึ้นสำหรับโมเดล) ไม่ใช่การลดความซับซ้อนของความสัมพันธ์โดยตรง

\subsection*{ข้อที่ 18}
\textbf{โจทย์:} โมเดลในตัวเลือกใดไม่จำเป็นต้องทำ Feature scaling\\
\textbf{คำตอบ:} Tree-based models\\
\textbf{คำอธิบาย:} การตัดสินใจของต้นไม้ใช้จุดตัด (Threshold) ทีละฟีเจอร์ ไม่ได้วัดระยะห่างระหว่างฟีเจอร์ จึงไม่ได้รับผลกระทบจากสเกล

\subsection*{ข้อที่ 19}
\textbf{โจทย์:} ข้อใดกล่าวถูกต้องเกี่ยวกับ SHAP\\
\textbf{คำตอบ:} เป็นวิธีที่ใช้ในการแปลความหมายของการทำนายว่า Features ใดมีความสำคัญเท่าไหร่\\
\textbf{คำอธิบาย:} SHAP ใช้ทฤษฎีเกมเพื่อจัดสรรส่วนแบ่งความสำคัญให้แต่ละ Feature ในการทำนายผลลัพธ์

\newpage
\section{Machine Learning Series: Unsupervised Learning}

\subsection*{ข้อที่ 1}
\textbf{โจทย์:} Supervised และ Unsupervised ต่างกันอย่างไร\\
\textbf{คำตอบ:} Supervised มีเฉลย (Label), Unsupervised ไม่มีเฉลย\\
\textbf{คำอธิบาย:} Unsupervised เรียนรู้จากโครงสร้างข้อมูลดิบเพื่อค้นหา Pattern โดยปราศจากป้ายกำกับ

\subsection*{ข้อที่ 2}
\textbf{โจทย์:} หลักการประเมิน Cluster ที่ดีคืออะไร\\
\textbf{คำตอบ:} ถูกทุกตัวเลือก (สมาชิกในกลุ่มเหมือนกัน, ต่างกลุ่มกันต่างกันมาก, จำนวนสมาชิกไม่น้อยเกินไป)\\
\textbf{คำอธิบาย:} Cluster ที่ดีควรมีความหนาแน่นภายในสูงและมีการแยกห่างระหว่างกลุ่มชัดเจนเพื่อให้เกิดนัยสำคัญ

\subsection*{ข้อที่ 3}
\textbf{โจทย์:} distance กับ similarity ต่างกันอย่างไร\\
\textbf{คำตอบ:} distance ยาวได้ไม่จำกัด แต่ similarity มีค่าสูงสุดจำกัด\\
\textbf{คำอธิบาย:} Distance วัดความห่าง (เริ่มที่ 0 ไปถึง $\infty$) ส่วน Similarity วัดความเหมือน (มักอยู่ในช่วง [0, 1] หรือ [-1, 1])

\subsection*{ข้อที่ 4}
\textbf{โจทย์:} เราจะประเมินผล clustering ได้อย่างไรบ้าง\\
\textbf{คำตอบ:} ถูกทุกตัวเลือก (ใช้ labeled data, วัดความต่างระหว่างกลุ่ม, วัดความแปรปรวนภายใน)\\
\textbf{คำอธิบาย:} การประเมินทำได้ทั้งแบบ External (เทียบ Label) และ Internal (ดูโครงสร้างทางคณิตศาสตร์)

\subsection*{ข้อที่ 5}
\textbf{โจทย์:} ตัวเลือกใดไม่ถูกต้องเกี่ยวกับระยะทางและความคล้าย\\
\textbf{คำตอบ:} Cosine similarity เหมาะกับข้อมูลที่แต่ละมิติมีการกระจายไม่เท่ากัน (เป็นข้อที่กำกวมที่สุด)\\
\textbf{คำอธิบาย:} Cosine Similarity เน้นวัดทิศทางเวกเตอร์ ไม่ใช่คุณสมบัติหลักในการแก้ปัญหาการกระจายข้อมูลที่ไม่เท่ากัน

\subsection*{ข้อที่ 6}
\textbf{โจทย์:} ตัวเลือกใดเป็นสาเหตุให้ต้องลดมิติของข้อมูล\\
\textbf{คำตอบ:} ถูกทุกตัวเลือก (แสดงผลได้ง่าย, ประหยัดพื้นที่, เลือกเฉพาะข้อมูลที่เป็นประโยชน์)\\
\textbf{คำอธิบาย:} เพื่อทำให้ข้อมูลอยู่ในรูปแบบที่เรียบง่ายขึ้น (Simplified Representation) โดยยังรักษาใจความสำคัญ

\subsection*{ข้อที่ 7}
\textbf{โจทย์:} จากรูปตัวเลือกใดเหมาะสมในการลดมิติข้อมูล Moons dataset ให้เหลือ 1 มิติ\\
\textbf{ลักษณะรูปภาพ:} ภาพแสดงชุดข้อมูล Moons dataset ซึ่งเป็นจุดข้อมูลที่เรียงตัวกันเป็นรูปพระจันทร์เสี้ยวสองวงที่ซ้อนเหลื่อมกันในระนาบ 2 มิติ\\
\textbf{คำตอบ:} t-SNE\\
\textbf{คำอธิบาย:} ข้อมูล Moons เป็น Non-linear การใช้ PCA จะทำให้ข้อมูลทับซ้อนกัน แต่ t-SNE สามารถรักษาโครงสร้างความใกล้เคียงได้ดีกว่า

\subsection*{ข้อที่ 8}
\textbf{โจทย์:} ผลลัพธ์ใดสอดคล้องกับพฤติกรรมของ K-means บนข้อมูล Moons dataset\\
\textbf{ลักษณะรูปภาพ:} ภาพแสดงผลลัพธ์การแบ่งกลุ่มข้อมูล Moons dataset ด้วยอัลกอริทึมที่แตกต่างกัน (A-D)\\
\textbf{คำตอบ:} รูปที่ D\\
\textbf{คำอธิบาย:} K-means พยายามแบ่งพื้นที่ด้วยเส้นตรง (Decision boundary) จึงได้ผลลัพธ์เป็นการตัดผ่านกลุ่มข้อมูลในแนวระนาบ

\subsection*{ข้อที่ 9}
\textbf{โจทย์:} เมื่อต้องการสร้าง cluster ของธุรกรรมออนไลน์ที่ประกอบด้วยยอดเงินและประเภทร้านค้า ควรเลือกตัวแทนอย่างไร\\
\textbf{คำตอบ:} Mean ของยอดเงิน, ช่วงเวลา และ mode ของ ประเภทร้านค้า, timezone\\
\textbf{คำอธิบาย:} ข้อมูลเชิงปริมาณใช้ค่าเฉลี่ย (Mean) ส่วนข้อมูลเชิงกลุ่มใช้ฐานนิยม (Mode) เป็นตัวแทนที่เหมาะสม

\subsection*{ข้อที่ 10}
\textbf{โจทย์:} ควรเลือกจำนวน cluster (k) อย่างไร\\
\textbf{คำตอบ:} ถูกทุกตัวเลือก (ทดลองทีละค่า, ตั้งมากกว่าที่คาดไว้เล็กน้อย, ใช้ elbow method)\\
\textbf{คำอธิบาย:} เป็นกระบวนการที่ใช้ร่วมกันเพื่อหาความสมดุลระหว่างความแม่นยำและความหมายทางธุรกิจ

\subsection*{ข้อที่ 11}
\textbf{โจทย์:} ตัวเลือกใดเป็นวิธีการแบ่งกลุ่มข้อมูล (cluster) แบบมีลำดับ\\
\textbf{คำตอบ:} แบ่งกลุ่มใหญ่ก่อนค่อยแบ่งย่อย และ สร้างกลุ่มย่อยก่อนค่อยรวมเป็นกลุ่มใหญ่\\
\textbf{คำอธิบาย:} ครอบคลุมทั้งวิธี Divisive (Top-down) และ Agglomerative (Bottom-up)

\subsection*{ข้อที่ 12}
\textbf{โจทย์:} ตัวเลือกใดเป็นการกำหนดค่าเพื่อให้ได้ cluster ที่เป็นประโยชน์\\
\textbf{คำตอบ:} กำหนดจำนวนสมาชิกขั้นต่ำใน cluster\\
\textbf{คำอธิบาย:} เพื่อป้องกันการสร้างกลุ่มที่เป็น Noise หรือกลุ่มที่มีสมาชิกน้อยเกินไปจนไม่มีนัยสำคัญทางสถิติ

\subsection*{ข้อที่ 13}
\textbf{โจทย์:} จุดเด่นของ density-based clustering คืออะไร\\
\textbf{คำตอบ:} รองรับรูปร่างของ cluster หลายรูปแบบ\\
\textbf{คำอธิบาย:} เนื่องจากทำงานตามความหนาแน่น ไม่ยึดติดกับจุดศูนย์กลาง จึงตรวจจับกลุ่มที่มีรูปร่างอิสระได้

\subsection*{ข้อที่ 14}
\textbf{โจทย์:} DBSCAN สร้าง cluster ได้อย่างไร\\
\textbf{คำตอบ:} รวม sample ในพื้นที่เล็ก ๆ ที่หนาแน่นพอไว้ด้วยกัน\\
\textbf{คำอธิบาย:} ใช้พารามิเตอร์ Eps และ MinPts ในการเชื่อมต่อจุดที่อยู่ในพื้นที่หนาแน่นสูงเข้าด้วยกัน

\subsection*{ข้อที่ 15}
\textbf{โจทย์:} ใน association rule เกณฑ์ใดกำหนดความน่าสนใจของ Frequent Itemset\\
\textbf{คำตอบ:} Support (ความถี่ในการปรากฏ)\\
\textbf{คำอธิบาย:} Support บอกว่ารูปแบบพฤติกรรมนั้นเกิดขึ้นบ่อยเพียงใดในภาพรวมของชุดข้อมูลทั้งหมด

\subsection*{ข้อที่ 16}
\textbf{โจทย์:} support กับ confidence ต่างกันอย่างไร\\
\textbf{คำตอบ:} support ใช้กับ itemset 1 ชุด ส่วน confidence ใช้บอกความสัมพันธ์ของ itemset 2 ชุด\\
\textbf{คำอธิบาย:} Support บอกความนิยมของสินค้าก้อนนั้น ส่วน Confidence บอกโอกาสที่สินค้า B จะถูกซื้อถ้าซื้อ A ไปแล้ว

\subsection*{ข้อที่ 17}
\textbf{โจทย์:} เมื่อทำ matrix factorization $A = BC$ ตัวเลือกใดถูกต้อง\\
\textbf{คำตอบ:} ถูกทุกตัวเลือก\\
\textbf{คำอธิบาย:} ช่วยลดพื้นที่จัดเก็บ, ดูการกระจายคำใน topic ได้จาก matrix C และแต่ละแถวคือ document vector

\subsection*{ข้อที่ 18}
\textbf{โจทย์:} เมื่อใช้ topic model จะหา topic vector ของเอกสารใหม่ที่ไม่ได้ใช้ train ได้อย่างไร\\
\textbf{คำตอบ:} กรณีที่เป็น LDA ให้ fit document ใหม่กับ model เดิม เพื่อหา topic distribution vector\\
\textbf{คำอธิบาย:} เป็นกระบวนการ Inference มาตรฐานของ LDA เพื่อหาพารามิเตอร์แฝงสำหรับข้อมูลใหม่

\subsection*{ข้อที่ 19}
\textbf{โจทย์:} ตัวเลือกใดเป็นตัวอย่างการใช้งาน semi-supervised learning\\
\textbf{คำตอบ:} เพิ่มความแม่นยำของ classification model โดยอาศัย unlabeled data\\
\textbf{คำอธิบาย:} ใช้ประโยชน์จากข้อมูลไม่มีฉลากปริมาณมากร่วมกับข้อมูลมีฉลากน้อยชิ้นเพื่อหาขอบเขตคลาสที่ดีขึ้น

\subsection*{ข้อที่ 20}
\textbf{โจทย์:} หากจัดกลุ่มแล้วพบว่าได้ 4 กลุ่มที่มีสมาชิกพอๆ กัน ข้อใดถูกต้อง\\
\textbf{คำตอบ:} ไม่สามารถสรุปได้\\
\textbf{คำอธิบาย:} อาจเป็นผลลัพธ์จากอัลกอริทึมที่บีบให้เท่ากัน หรือเป็นเรื่องบังเอิญ ไม่ได้การันตีความเหมาะสม

\subsection*{ข้อที่ 21}
\textbf{โจทย์:} หากต้องเลือกระหว่างวิธี A B และ C ในการจัดกลุ่ม ควรใช้เกณฑ์ใด\\
\textbf{คำตอบ:} ถูกทุกตัวเลือก (เทียบ variance, เทียบ similarity, เทียบ labeled dataset)\\
\textbf{คำอธิบาย:} เป็นเกณฑ์มาตรฐานในการประเมินประสิทธิภาพทั้งแบบเชิงโครงสร้างและแบบอ้างอิงความจริง

\newpage
\section{Machine Learning Series: Time Series Modeling}

\subsection*{ข้อที่ 1}
\textbf{โจทย์:} ข้อใดไม่ใช่ข้อมูลอนุกรมเวลา (Time Series)\\
\textbf{คำตอบ:} ค่าไฟฟ้าของห้องทุกห้องในคอนโดตอนเดือนมกราคม 2020\\
\textbf{คำอธิบาย:} เป็นข้อมูลตัดขวาง (Cross-sectional) ณ จุดเวลาเดียว ไม่ได้ดูการเปลี่ยนแปลงตามลำดับเวลาของหน่วยเดิม

\subsection*{ข้อที่ 2}
\textbf{โจทย์:} ข้อใดไม่ถูกต้องเกี่ยวกับข้อมูล time series\\
\textbf{คำตอบ:} ข้อมูล time series ต้องเก็บนนช่วงเวลาที่ห่างเท่า ๆ กัน\\
\textbf{คำอธิบาย:} ไม่จำเป็นเสมอไป ข้อมูลบางประเภทเช่นราคาหุ้นรายธุรกรรมอาจบันทึกในช่วงเวลาที่ไม่สม่ำเสมอได้

\subsection*{ข้อที่ 3}
\textbf{โจทย์:} ข้อมูลจำนวนผู้ป่วยใหม่รายวันในประเทศต่างๆ ดังรูป เป็นข้อมูลประเภทใด\\
\textbf{ลักษณะรูปภาพ:} ภาพแสดงกราฟเส้นจำนวนผู้ป่วย COVID-19 หลายเส้นตามแต่ละประเทศในช่วงเวลาเดียวกัน\\
\textbf{คำตอบ:} เป็น time series แบบ multivariate\\
\textbf{คำอธิบาย:} เนื่องจากมีการพิจารณาหลายตัวแปร (หลายประเทศ) ที่เปลี่ยนแปลงไปตามลำดับเวลาเดียวกัน

\subsection*{ข้อที่ 4}
\textbf{โจทย์:} ลักษณะใดที่บ่งบอกว่าข้อมูลเป็น Additive time series\\
\textbf{ลักษณะรูปภาพ:} ภาพแสดงกราฟความต้องการใช้ไฟฟ้า (Electricity Demand) ที่มีความผันผวนคงที่ตลอดช่วงเวลา\\
\textbf{คำตอบ:} Additive\\
\textbf{คำอธิบาย:} ใน Additive model ขนาดของความผันผวน (Seasonal variation) จะคงที่ ไม่แปรผันตามระดับแนวโน้มหลัก

\subsection*{ข้อที่ 5}
\textbf{โจทย์:} จากกราฟ A, B, C จงระบุส่วนประกอบจากการ Decompose\\
\textbf{ลักษณะรูปภาพ:} ภาพแสดงองค์ประกอบ 3 ส่วน ได้แก่ กราฟ A (คลื่นซ้ำๆ), กราฟ B (สุ่ม), กราฟ C (เส้นโค้งทิศทาง)\\
\textbf{คำตอบ:} A – seasonal, B – residual, C – trend\\
\textbf{คำอธิบาย:} Seasonal แสดงรูปแบบซ้ำๆ, Trend แสดงทิศทางระยะยาว และ Residual แสดงค่าที่เหลือที่เป็นสัญญาณรบกวน

\subsection*{ข้อที่ 6}
\textbf{โจทย์:} ข้อใดไม่ถูกต้องเกี่ยวกับ time series decomposition\\
\textbf{คำตอบ:} ข้อมูล time series จะประกอบด้วยสามส่วนเสมอ (และ Trend บ่งบอกความสัมพันธ์ปัจจุบันกับอดีต)\\
\textbf{คำอธิบาย:} ในทางวิชาการอาจมีส่วนประกอบอื่นเช่น Cyclical และ Trend บอกแค่ทิศทาง ไม่ได้บอกความสัมพันธ์เชิงโครงสร้างย้อนหลัง

\subsection*{ข้อที่ 7}
\textbf{โจทย์:} ข้อใดกล่าวถูกต้องที่สุดเกี่ยวกับ stationary property ของ time series\\
\textbf{คำตอบ:} ถ้าแบ่งข้อมูล stationary เป็นสองส่วน ค่า variance ของครึ่งแรกจะมีค่าเท่ากับของครึ่งหลัง\\
\textbf{คำอธิบาย:} ตามนิยามของ Stationary ความแปรปรวน (Variance) ต้องคงที่ตลอดช่วงเวลา

\subsection*{ข้อที่ 8}
\textbf{โจทย์:} หลักฐานใดใช้สนับสนุนว่าข้อมูลเป็น stationary time series\\
\textbf{คำตอบ:} A, C, D, F (Running mean คงที่, ไม่มี seasonal, ADF p-value ต่ำ, Autocorrelation ลดลงเร็ว)\\
\textbf{คำอธิบาย:} เป็นคุณสมบัติพื้นฐานและผลการทดสอบทางสถิติที่ยืนยันความคงที่ของข้อมูล

\subsection*{ข้อที่ 10}
\textbf{โจทย์:} ถ้าจะแปลงข้อมูลให้กลายเป็น stationary จะต้องทำกระบวนการใดบ้าง\\
\textbf{คำตอบ:} A, B, C (ทำ Differencing ลด trend/seasonal และทำ log transformation)\\
\textbf{คำอธิบาย:} เพื่อจัดการทั้งความแปรปรวนที่ไม่คงที่, อิทธิพลของฤดูกาล และแนวโน้มของข้อมูล

\subsection*{ข้อที่ 11}
\textbf{โจทย์:} ข้อใดไม่ใช่ประโยชน์ของการ smoothing\\
\textbf{คำตอบ:} สามารถนำมาใช้ทำนายอนาคตได้\\
\textbf{คำอธิบาย:} Smoothing เป็นการจัดการข้อมูลอดีตเพื่อให้มองเห็น Pattern ชัดขึ้น แต่ไม่มีคุณสมบัติในการพยากรณ์อนาคตในตัวมันเอง

\subsection*{ข้อที่ 12}
\textbf{โจทย์:} หากข้อมูลมี Seasonality แต่ไม่มี Trend ควรใช้ smoothing ประเภทใดดีที่สุด\\
\textbf{คำตอบ:} Triple Exponential Smoothing (Holt-Winters)\\
\textbf{คำอธิบาย:} แม้ระบุว่าไม่มี Trend แต่ในตระกูลนี้ มีเพียง Triple Smoothing เท่านั้นที่รองรับองค์ประกอบฤดูกาล (Seasonality)

\subsection*{ข้อที่ 13}
\textbf{โจทย์:} ข้อใดเป็นเหตุผลหลักที่จะใช้ moving average แทน exponential moving average\\
\textbf{คำตอบ:} เมื่อต้องการให้น้ำหนักของข้อมูลในการเฉลี่ยเท่า ๆ กัน\\
\textbf{คำอธิบาย:} Simple Moving Average ให้น้ำหนักทุกจุดเท่ากัน ต่างจาก EMA ที่ให้น้ำหนักข้อมูลใหม่มากกว่า

\subsection*{ข้อที่ 14}
\textbf{โจทย์:} SARIMA, ARIMA, ARMA จะทำงานได้ดีบนข้อมูลลักษณะใด\\
\textbf{คำตอบ:} ข้อมูล Stationary และ ลักษณะของข้อมูลในอดีตบ่งบอกถึงลักษณะข้อมูลในอนาคต\\
\textbf{คำอธิบาย:} ต้องมีความคงที่ทางสถิติและมีความสัมพันธ์กันเอง (Autocorrelation) เพื่อให้โมเดลเรียนรู้ได้

\subsection*{ข้อที่ 15}
\textbf{โจทย์:} ข้อใดไม่ถูกต้องเกี่ยวกับทำนายบนข้อมูล time series\\
\textbf{คำตอบ:} โมเดล SARIMA จะได้ผลดีที่สุด เพราะสามารถรับมือข้อมูลได้ทุกประเภท\\
\textbf{คำอธิบาย:} ไม่มีโมเดลเดียวที่เก่งที่สุด SARIMA เป็นเชิงเส้น หากข้อมูลซับซ้อนมาก Machine Learning อาจทำได้ดีกว่า

\subsection*{ข้อที่ 16}
\textbf{โจทย์:} ข้อใดไม่ถูกต้องเกี่ยวกับโมเดล ARMA, ARIMA และ SARIMA\\
\textbf{คำตอบ:} เวลาใช้โมเดลให้เลือกค่า p q ที่ทำให้ได้ค่า error น้อยที่สุดใน training data\\
\textbf{คำอธิบาย:} การเลือกตาม training error อย่างเดียวจะนำไปสู่ Overfitting ควรใช้เกณฑ์เช่น AIC/BIC

\subsection*{ข้อที่ 17}
\textbf{โจทย์:} ข้อใดถูกต้องที่สุดเกี่ยวกับ confidence interval\\
\textbf{คำตอบ:} Confidence interval จะกว้างขึ้นตาม % ค่า confidence\\
\textbf{คำอธิบาย:} เพื่อความมั่นใจที่สูงขึ้น ต้องขยายช่วงให้กว้างขึ้นเพื่อดักจับค่าที่แท้จริงได้ครอบคลุมมากขึ้น

\subsection*{ข้อที่ 18}
\textbf{โจทย์:} ข้อใดไม่ถูกต้องเกี่ยวกับ Facebook Prophet\\
\textbf{คำตอบ:} สามารถทำทุกอย่างที่ SARIMA ทำได้\\
\textbf{คำอธิบาย:} สองโมเดลมีปรัชญาต่างกัน Prophet เป็น Additive Curve Fitting ส่วน SARIMA เป็น Stochastic Process

\subsection*{ข้อที่ 19}
\textbf{โจทย์:} กรณีใดที่ควรใช้ supervised learningแทน time series modeling\\
\textbf{คำตอบ:} มีข้อมูลน้อย เช่น มีรอบ seasonalเพียงหนึ่งรอบ\\
\textbf{คำอธิบาย:} โมเดล Time Series ต้องการข้อมูลหลายรอบเพื่อแยก Trend/Seasonal หากน้อยไปควรใช้ Feature Engineering ใน Supervised แทน

\subsection*{ข้อที่ 20}
\textbf{โจทย์:} ข้อใดไม่ถูกต้องเกี่ยวกับการโมเดล time series\\
\textbf{คำตอบ:} ควรใช้โมเดลที่ซับซ้อนที่สุดเพื่อจะได้รับมือกับข้อมูลทุกรูปแบบ\\
\textbf{คำอธิบาย:} ขัดกับหลัก Occam's Razor โมเดลซับซ้อนเกินไปมัก Overfit และอธิบายได้ยาก

\subsection*{ข้อที่ 21}
\textbf{โจทย์:} ข้อใดไม่ถูกต้องเกี่ยวกับการประยุกต์ใช้ time series modeling\\
\textbf{คำตอบ:} สามารถนำ Confidence interval ของการทำนายใช้บอกจุดสูงสุดและต่ำสุดของความเป็นไปได้ของข้อมูลได้\\
\textbf{คำอธิบาย:} ต้องใช้ Prediction Interval (PI) ไม่ใช่ Confidence Interval (CI) เนื่องจาก PI ครอบคลุมความผันผวนของข้อมูลจริงด้วย

\end{document}
